\fintraceabstractheader
\begin{fintraceabstractpanel}
\begin{abstract}

大语言模型为上市公司研究提供了新的自动化路径，但现有金融智能体仍面临复杂任务规划成本高、多源异构数据难以协同、长周期执行不稳定及生成结论缺乏证据约束等问题。本文设计并实现 \textbf{FinTrace：}面向上市公司研究的证据驱动可追溯金融智能体，将金融研究任务转化为可规划、可执行、可校验和可追溯的结构化调查过程。围绕“任务规划--受控执行--多源取证--可信生成”的技术链路，FinTrace 形成以下主要创新：

\begin{itemize}[leftmargin=6mm,label=\textcolor{FinOrange}{\raisebox{.12ex}{\large$\bullet$}},itemsep=5pt,topsep=5pt]

\item \fintraceabstractitemtitle{分层任务规划：确定性直连与受限 ReAct 相结合的规划机制：}
解决“如何规划”的问题。对于任务明确、参数充分且动作唯一的请求，直接执行确定性工具调用；对于复杂研究任务，则由 Planner 按“规划--执行--观察--再规划”逐步推进，并通过步骤、工具调用及新增证据约束调查范围，在保留推理能力的同时降低冗余规划与失控循环风险。

\item \fintraceabstractitemtitle{受控任务执行：面向金融研究的 Agent Harness：}
解决“如何稳定执行”的问题。系统采用“主 Agent + 专业确定性工具”架构，基于 LangGraph 组织请求解析、动作校验、工具执行、结果验证、证据合并与动作修复，并结合 Pydantic/Schema 对实体、时间、参数和返回结果实施结构化约束，确保复杂任务执行稳定可审计。

\item \fintraceabstractitemtitle{时序证据组织：多源异构金融数据协同方法：}
解决“如何获取可信证据”的问题。统一组织财务报表、股权快照、公司公告、机构研报及事件序列等异构数据，将结构化事实与可定位文本 Chunk 协同接入专业分析工具，并统一报告期、发布日期、快照日期和研究截止时点等时间语义，实现多源信息的组合调用与交叉验证。

\item \fintraceabstractitemtitle{可信结论生成：证据驱动的研究--验证--生成闭环：}
解决“如何形成可信结论”的问题。系统将 Evidence 作为工具执行与最终回答之间的核心数据对象，在执行结束后汇总调查结果并进行证据充分性审查：证据不足时重新触发 Planner 补充取证；通过 Evidence ID 建立结论与来源之间的对应关系，实现研究过程与结论的可追溯。

\end{itemize}

FinTrace从数据处理、索引构建、任务规划与工具执行到证据审查、结果生成和在线交互的端到端实现，并围绕财务、股权、事件、研报及综合研究任务建立系统评测流程。系统部署于 \url{https://fintrace.huayingkai.cn}，可直接在线体验（用户名：fintrace；密码：fintrace2026）；网站同时展示基于1410 条测试数据及其对应的系统回答，供查阅与复核。

\keywords{金融智能体，上市公司研究，Agent Harness，证据驱动，可追溯生成，Agentic AI}

\end{abstract}
\end{fintraceabstractpanel}
